%% toc: P3 32 32.1,32.2,32.3,32.4
\section{Proofs for stability and the forking calculus}
\label{sec:annex-p03-w03-p03-031-forking-calculus}

Throughout, \(T\), \(\bar\kappa\) and \(\mathfrak C\) are as in
\cref{def:p03-w03-p03-031-monster}, and all sets and tuples are small.  A
set of formulas is \emph{inconsistent} when it is not satisfied in
\(\mathfrak C\); by compactness and elementarity this is the same as
inconsistency with the elementary diagram of any model containing the
parameters.  If \(\{\varphi(x;b_i):i\in I\}\) is inconsistent, some finite
\(F\subseteq I\) has \(\mathfrak C\entails\neg\exists
x\bigwedge_{i\in F}\varphi(x;b_i)\); this is a formula true of
\((b_i)_{i\in F}\), and we use repeatedly that it is therefore true of any
tuple with the same type over the parameters of \(\varphi\).  For an
\(A\)-indiscernible sequence it then holds for every increasing tuple of
the same length.  We write \(Ab\) for \(A\cup\{\text{entries of }b\}\).

\begin{annexitem}{annex:p03-w03-p03-031-cb}{thm:p03-w03-p03-031-cb}{Cantor--Bendixson ranks of types}
\begin{proof}
Write \(X=S_m(A)\) and \(\operatorname{CB}\) for \(\operatorname{CB}_A\).
Every \(p\in X\) is a set of formulas in \(m\) variables with parameters
from the small set \(A\), each finite subset of which is satisfied, so it is
realized in \(\mathfrak C\) by \cref{thm:p03-w03-p03-031-saturation}~(1).
If \(c\) realizes \(p\), then \(p=\operatorname{tp}(c/A)\); hence \(p\)
contains \(\theta\wedge\chi\) whenever it contains \(\theta\) and \(\chi\),
it contains exactly one of \(\chi\) and \(\neg\chi\) for every formula
\(\chi\) in \(x\) over \(A\), and it contains \(\theta'\) whenever it
contains \(\theta\) and \(\mathfrak C\entails\forall x\,(\theta\to\theta')\).

\emph{Covering.}  We use the following fact.
\begin{itemize}
  \item[\((\ast)\)] Let \(\theta\) be a formula in \(x\) over \(A\), and for
        each \(p\in[\theta]\) let \(\chi_p\in p\) be given.  Then there are
        \(p_1,\ldots,p_n\in[\theta]\) with
        \([\theta]\subseteq[\chi_{p_1}]\cup\cdots\cup[\chi_{p_n}]\).
\end{itemize}
The set \(\Sigma=\{\theta\}\cup\{\neg\chi_p:p\in[\theta]\}\) is not realized
in \(\mathfrak C\): if \(c\) realized it, then
\(q=\operatorname{tp}(c/A)\) would lie in \([\theta]\), and \(c\) would
satisfy \(\neg\chi_q\) although \(\chi_q\in q\).  The parameters of
\(\Sigma\) lie in \(A\) and it has \(m\) variables, so by
\cref{thm:p03-w03-p03-031-saturation}~(1) some finite subset of \(\Sigma\)
is not satisfied in \(\mathfrak C\).  Adding \(\theta\) to it, we find
\(p_1,\ldots,p_n\in[\theta]\) such that
\(\theta\wedge\neg\chi_{p_1}\wedge\cdots\wedge\neg\chi_{p_n}\) is not
satisfied.  If \(q\in[\theta]\) and \(c\) realizes \(q\), then \(c\)
satisfies \(\theta\), so it fails some \(\neg\chi_{p_i}\), and
\(\chi_{p_i}\in q\).

\emph{Item 1.}  By clause~2 of \cref{def:p03-w03-p03-031-cb-rank},
\(X^{(\alpha+1)}\subseteq X^{(\alpha)}\).  By induction on \(\beta\),
\(X^{(\beta)}\subseteq X^{(\alpha)}\) for all \(\alpha\leq\beta\): the case
\(\beta=\alpha\) is trivial; at a successor \(\beta=\gamma+1>\alpha\) we
have \(\gamma\geq\alpha\) and \(X^{(\gamma+1)}\subseteq
X^{(\gamma)}\subseteq X^{(\alpha)}\); and at a limit \(\beta>\alpha\),
\(X^{(\beta)}=\bigcap_{\gamma<\beta}X^{(\gamma)}\subseteq X^{(\alpha)}\)
because \(\alpha<\beta\).  Hence the sets
\(X^{(\alpha)}\setminus X^{(\alpha+1)}\) are pairwise disjoint, and a type
lies in at most one of them.

Let \(\kappa=\abs X^+\).  If \(X^{(\alpha+1)}\neq X^{(\alpha)}\) for every
\(\alpha<\kappa\), choose \(p_\alpha\in X^{(\alpha)}\setminus
X^{(\alpha+1)}\).  For \(\alpha<\beta<\kappa\),
\(p_\beta\in X^{(\beta)}\subseteq X^{(\alpha+1)}\), so
\(p_\beta\neq p_\alpha\); this is an injection of \(\kappa\) into \(X\),
which is impossible.  So some \(\alpha^*\) has
\(X^{(\alpha^*+1)}=X^{(\alpha^*)}\).  By induction on \(\beta\geq\alpha^*\),
\(X^{(\beta)}=X^{(\alpha^*)}\): if \(X^{(\beta)}=X^{(\alpha^*)}\), then
\(X^{(\beta+1)}\) is obtained from \(X^{(\beta)}\) by the rule that
obtains \(X^{(\alpha^*+1)}\) from \(X^{(\alpha^*)}\), so it equals
\(X^{(\alpha^*+1)}=X^{(\alpha^*)}\); at a limit \(\delta>\alpha^*\),
\(X^{(\delta)}\) is the intersection of a decreasing family that is
constantly \(X^{(\alpha^*)}\) from \(\alpha^*\) on.

If \(p\in X^{(\alpha^*)}\), then \(p\in X^{(\beta)}\) for every \(\beta\),
by the two facts just proved, so \(\operatorname{CB}(p)=\infty\).  If
\(p\notin X^{(\alpha^*)}\), let \(\gamma\) be least with
\(p\notin X^{(\gamma)}\).  Then \(\gamma\neq0\) because \(X^{(0)}=X\), and
\(\gamma\) is not a limit, since otherwise \(p\notin X^{(\alpha)}\) for some
\(\alpha<\gamma\).  So \(\gamma=\alpha+1\) with \(p\in X^{(\alpha)}\), and
\(\operatorname{CB}(p)=\alpha\).  By the disjointness shown above, the rank
is unique.

\emph{Item 2.}  Let \(\operatorname{CB}(p)=\alpha\).  Then
\(p\in X^{(\alpha)}\) and \(p\notin X^{(\alpha+1)}\), so by clause~2 of
the definition some \(\theta\in p\) is such that \([\theta]\cap
X^{(\alpha)}\) contains no type other than \(p\).  If \(\beta>\alpha\),
then \(X^{(\beta)}\subseteq X^{(\alpha+1)}\subseteq X^{(\alpha)}\) and
\(p\notin X^{(\alpha+1)}\), so \([\theta]\cap X^{(\beta)}\subseteq
\{p\}\cap X^{(\beta)}=\varnothing\).

\emph{Item 3.}  Let \(Y=X^{(\alpha^*)}\neq\varnothing\), and call a
formula \(\theta\) in \(x\) over \(A\) \emph{large} when
\([\theta]\cap Y\neq\varnothing\).  A large formula belongs to a type and
is therefore realized.  We claim: if \(\theta\) is large, some formula
\(\chi\) in \(x\) over \(A\) makes both \(\theta\wedge\chi\) and
\(\theta\wedge\neg\chi\) large.  Take \(p\in[\theta]\cap Y\).  As
\(Y=X^{(\alpha^*+1)}\) and \(\theta\in p\), the set
\([\theta]\cap X^{(\alpha^*)}=[\theta]\cap Y\) contains a type
\(q\neq p\).  Two distinct complete types over \(A\) differ on some
formula, so some \(\chi\) has \(\chi\in p\) and \(\neg\chi\in q\).  Then
\(\theta\wedge\chi\in p\) and \(\theta\wedge\neg\chi\in q\), and
\(p,q\in Y\).

The formula \(x=x\) belongs to every type, so it is large.  Choose formulas
\(\theta_s\), for \(s\) in the set \(2^{<\omega}\) of finite binary
sequences, by recursion on the length of \(s\): \(\theta_\varnothing\) is
\(x=x\); if \(\theta_s\) is large, take \(\chi_s\) as in the claim and put
\(\theta_{s\frown\langle0\rangle}=\theta_s\wedge\chi_s\) and
\(\theta_{s\frown\langle1\rangle}=\theta_s\wedge\neg\chi_s\), both large.
For \(\sigma\in2^\omega\), every finite subset of
\(\{\theta_{\sigma\mathbin{\upharpoonright}n}:n<\omega\}\) is implied by its
member of greatest length, which is large and so realized; by
\cref{thm:p03-w03-p03-031-saturation}~(1) the whole set is realized by some
\(c_\sigma\).  Put \(p_\sigma=\operatorname{tp}(c_\sigma/A)\in X\).  If
\(\sigma\neq\tau\), let \(n\) be least with \(\sigma(n)\neq\tau(n)\) and
\(s=\sigma\mathbin{\upharpoonright}n=\tau\mathbin{\upharpoonright}n\); one of
\(p_\sigma\), \(p_\tau\) contains \(\theta_s\wedge\chi_s\) and the other
contains \(\theta_s\wedge\neg\chi_s\), so \(p_\sigma\neq p_\tau\).  Thus
\(\sigma\mapsto p_\sigma\) is injective and \(\abs X\geq2^{\aleph_0}\).
If \(X\) is countable, then \(\abs X<2^{\aleph_0}\) by Cantor's theorem, so
\(X^{(\alpha^*)}=\varnothing\), and by item~1 every type has an ordinal
rank.

\emph{Item 4.}  The formula \(\theta\) is realized, so \([\theta]\neq
\varnothing\).  For each \(p\in[\theta]\), item~2 gives \(\theta_p\in p\)
with \([\theta_p]\cap X^{(\operatorname{CB}(p))}\subseteq\{p\}\).  By
\((\ast)\) there are \(p_1,\ldots,p_n\in[\theta]\), with \(n\geq1\), such
that every \(q\in[\theta]\) lies in some \([\theta_{p_i}]\).  Let
\(\alpha_\theta=\max_i\operatorname{CB}(p_i)\).

Let \(q\in[\theta]\) and choose \(i\) with \(q\in[\theta_{p_i}]\).  If
\(\operatorname{CB}(q)>\operatorname{CB}(p_i)\), then \(q\in
X^{(\operatorname{CB}(q))}\subseteq X^{(\operatorname{CB}(p_i))}\), so
\(q\in[\theta_{p_i}]\cap X^{(\operatorname{CB}(p_i))}\subseteq\{p_i\}\) and
\(q=p_i\), which contradicts \(\operatorname{CB}(q)>\operatorname{CB}(p_i)\).
Hence \(\operatorname{CB}(q)\leq\operatorname{CB}(p_i)\leq\alpha_\theta\),
and \(\alpha_\theta\), which is the rank of one of the \(p_i\), is the
greatest rank of a type in \([\theta]\).  If moreover
\(\operatorname{CB}(q)=\alpha_\theta\), then
\(\operatorname{CB}(p_i)=\alpha_\theta\) and \(q\in[\theta_{p_i}]\cap
X^{(\alpha_\theta)}\subseteq\{p_i\}\).  So every type of rank
\(\alpha_\theta\) in \([\theta]\) is one of \(p_1,\ldots,p_n\).
\end{proof}
\end{annexitem}

\begin{annexitem}{annex:p03-w03-p03-031-omega-ss}{thm:p03-w03-p03-031-omega-superstable}{\(\omega\)-stable theories are superstable}
\begin{proof}
The saturation cardinal satisfies \(\bar\kappa>2^{\abs T}\geq\abs T\geq
\aleph_0\), so every countable set is small.

\emph{Item 1.}  By \cref{def:p03-stability-spectrum}, \(\omega\)-stability
is \(\aleph_0\)-stability, so \(T\) is \(\lambda\)-stable for the infinite
cardinal \(\lambda=\aleph_0\).  By the implication from~(2) to~(1) of
\cref{thm:p03-w03-p03-032-spectrum}, \(T\) is stable.

\emph{Item 2} is \cref{thm:p03-w03-p03-032-one-to-many} with
\(\lambda=\aleph_0\).

\emph{Item 3.}  Suppose \(\operatorname{Deg}_m(x=x)=\infty\) for some
\(m\geq1\).  Apply \cref{thm:p03-w03-p03-032-degree-tree} with the
infinite small cardinal \(\lambda=\aleph_0\).  It gives formulas
\(\psi_n(x;y_n)\) and integers \(k_n\geq2\), for \(n\geq1\), and tuples
\(a_\eta\), for \(\eta\in\omega^{<\omega}\setminus\{\varnothing\}\), such
that (i) for every \(\nu\in\omega^\omega\) the set
\(\{\psi_n(x;a_{\nu\mathbin{\upharpoonright}n}):n\geq1\}\) is realized, and
(ii) for every \(n<\omega\) and \(\eta\in\omega^n\) the family
\((\psi_{n+1}(x;a_{\eta\frown\langle i\rangle}):i<\omega)\) is
\(k_{n+1}\)-inconsistent.  Let \(A\) be the set of entries of the tuples
\(a_\eta\).  There are countably many nodes and each tuple is finite, so
\(A\) is countable, and \(X=S_m(A)\) is countable by item~2.  By
\cref{thm:p03-w03-p03-031-cb}~(3) every type in \(X\) has an ordinal rank
\(\operatorname{CB}_A\).

For \(\eta\in\omega^{<\omega}\) let
\[
 \theta_\eta\defeq(x=x)\wedge\bigwedge_{1\leq j\leq\abs\eta}
 \psi_j(x;a_{\eta\mathbin{\upharpoonright}j}),
\]
a formula over \(A\).  It is realized: if \(\nu\in\omega^\omega\) extends
\(\eta\), a realization of the set in (i) satisfies each conjunct.  By
\cref{thm:p03-w03-p03-031-cb}~(4), \([\theta_\eta]\) has a greatest rank
\(\alpha_\eta\), and the set \(F_\eta\) of types in \([\theta_\eta]\) of
rank \(\alpha_\eta\) is finite.

We claim that for every \(\eta\) with \(\abs\eta=n\) there is \(i<\omega\)
with \(\alpha_{\eta\frown\langle i\rangle}<\alpha_\eta\).  For
\(p\in F_\eta\), let \(I_p\) be the set of \(i<\omega\) with
\(\psi_{n+1}(x;a_{\eta\frown\langle i\rangle})\in p\).  Then
\(\abs{I_p}<k_{n+1}\): otherwise a realization of \(p\) would satisfy
\(k_{n+1}\) members of the \(k_{n+1}\)-inconsistent family of (ii).  The set
\(I=\bigcup_{p\in F_\eta}I_p\) is therefore finite; choose \(i\notin I\).
Let \(q\in[\theta_{\eta\frown\langle i\rangle}]\).  Since
\(\theta_{\eta\frown\langle i\rangle}\) is
\(\theta_\eta\wedge\psi_{n+1}(x;a_{\eta\frown\langle i\rangle})\), the type
\(q\) contains \(\theta_\eta\) and
\(\psi_{n+1}(x;a_{\eta\frown\langle i\rangle})\).  So \(q\in[\theta_\eta]\)
and \(\operatorname{CB}_A(q)\leq\alpha_\eta\); and \(q\notin F_\eta\), since
otherwise \(i\in I_q\subseteq I\).  As \(F_\eta\) contains every type of
\([\theta_\eta]\) of rank \(\alpha_\eta\),
\(\operatorname{CB}_A(q)<\alpha_\eta\).  This holds for every
\(q\in[\theta_{\eta\frown\langle i\rangle}]\), so
\(\alpha_{\eta\frown\langle i\rangle}<\alpha_\eta\).

Now define \(\eta_0=\varnothing\) and
\(\eta_{n+1}=\eta_n\frown\langle i_n\rangle\), where \(i_n\) is given by the
claim for \(\eta_n\).  Then
\(\alpha_{\eta_0}>\alpha_{\eta_1}>\alpha_{\eta_2}>\cdots\) is an infinite
strictly decreasing sequence of ordinals, which does not exist because the
ordinals are well ordered.  Hence \(\operatorname{Deg}_m(x=x)<\infty\) for
every \(m\geq1\).

\emph{Item 4.}  By items~1 and~3, \(T\) satisfies condition~(3) of
\cref{thm:p03-w03-p03-032-superstable}, which implies condition~(4):
\(T\) is superstable in the sense of \cref{def:p03-stability-spectrum}.
For the precise form, let \(p\) be a complete type in \(m\) variables over a
small set \(B\).  Every finite \(p_0\subseteq p\) has
\(\operatorname{Deg}\bigwedge p_0\leq\operatorname{Deg}_m(x=x)<\infty\) by
\cref{thm:p03-w03-p03-032-degree-basics}~(1), since \(\bigwedge p_0\)
implies \(x=x\).  These degrees are ordinals, so their least value
\(\operatorname{Deg}p\) is attained at some finite \(q\subseteq p\).  The
final clause of \cref{thm:p03-w03-p03-032-superstable} then says that \(p\)
does not fork over the finite set of parameters of \(q\).

\emph{A second proof of item~4 for a countable language.}  If the language
of \(T\) is countable, \(\omega\)-stability says that there are countably
many complete \(1\)-types over every countable parameter set in a model, so
by \cref{thm:p03-morley-countability} \(T\) is totally transcendental and
\(\abs{S_1(A)}\leq\abs A+\aleph_0\) for every parameter set \(A\) in a
model.  Hence \(T\) is \(\lambda\)-stable for every infinite \(\lambda\),
in particular for every \(\lambda\geq2^{\abs T}\).  This is condition~(1)
of \cref{thm:p03-w03-p03-032-superstable}, which implies condition~(4).
\end{proof}
\end{annexitem}

\begin{annexitem}{annex:p03-w03-p03-031-sat}{thm:p03-w03-p03-031-saturation}{Small types and witnesses inside the ambient model}
\begin{proof}
\emph{Item 1.}  Write \(q=q(x_\alpha:\alpha<\mu)\) with single variables
\(x_\alpha\).  We choose \(c_\alpha\in\mathfrak C\) by recursion on
\(\alpha\leq\mu\), keeping the invariant that the set \(q_\alpha\),
obtained from \(q\) by substituting \(c_\beta\) for \(x_\beta\) for all
\(\beta<\alpha\), is finitely satisfiable in \(\mathfrak C\).  For
\(\alpha=0\) this is the hypothesis.  At a limit \(\alpha\), a finite
subset of \(q_\alpha\) comes from finitely many formulas of \(q\) and
involves finitely many substituted \(c_\beta\), all with \(\beta<\gamma\)
for some \(\gamma<\alpha\); it is a finite subset of \(q_\gamma\) with some
further substitutions of the elements \(c_\beta\), \(\gamma\leq\beta\),
none of which occur, so it is satisfiable by the invariant at \(\gamma\).
At a successor step, given \(q_\alpha\), let
\[
 r(x_\alpha)=\Bigl\{\exists\bar y\,\bigwedge\Phi:\Phi\subseteq q_\alpha
   \text{ finite},\ \bar y\text{ the variables of }\Phi\text{ other than }
   x_\alpha\Bigr\}.
\]
Its parameters lie in \(D\cup\{c_\beta:\beta<\alpha\}\), a set of
cardinality at most \(\abs D+\abs\alpha<\bar\kappa\).  A finite subset of
\(r\) is implied by the single member built from the union of the finitely
many sets \(\Phi\) involved, which is satisfiable by the invariant; so
\(r\) is finitely satisfiable in \(\mathfrak C\).  By compactness \(r\) is
realized in an elementary extension of \(\mathfrak C\), and the complete
type of such a realization over the parameter set is consistent with
\(\mathfrak C\); by \(\bar\kappa\)-saturation
it is realized by some \(c_\alpha\in\mathfrak C\).  A finite subset of
\(q_{\alpha+1}\) arises from a finite \(\Phi\subseteq q_\alpha\) by
substituting \(c_\alpha\) for \(x_\alpha\), and \(c_\alpha\) satisfies
\(\exists\bar y\bigwedge\Phi\), so the invariant holds at \(\alpha+1\).  At
\(\alpha=\mu\) no variable remains, and each member of \(q_\mu\) is a
satisfiable, hence true, sentence; so \((c_\alpha)_{\alpha<\mu}\)
realizes \(q\).

\emph{Item 2.}  Let \(q(y)\) be the set of formulas \(\psi(c',y,d)\) such
that \(d\) is from \(A\) and \(\psi(c,e,d)\) holds in \(\mathfrak C\).  If
\(\psi_1(c',y,d_1),\ldots,\psi_m(c',y,d_m)\in q\), then
\(\exists y\bigwedge_k\psi_k(x,y,d_k)\) is a formula over \(A\) true of
\(c\), hence of \(c'\).  So \(q\) is finitely satisfiable, and item~1
gives a realization \(e'\).  Every formula over \(A\) true of \(ce\) is then
true of \(c'e'\); since both types are complete, \(ce\equiv_Ac'e'\).

\emph{Item 3.}  First let the witnesses lie in an elementary extension
\(\mathcal N\) of \(\mathfrak C\).  Suppose
\((b_i)_{i<\omega}\) in \(\mathcal N\) is \(A\)-indiscernible with
\(b_0=b\) and \(\{\varphi(x;b_i)\}\) inconsistent.  The set
\(\operatorname{tp}^{\mathcal N}((b_i)_{i\geq1}/Ab)\) has parameters in
\(\mathfrak C\), and each finite part is satisfied in \(\mathcal N\), so
the corresponding existential sentence holds in \(\mathcal N\) and hence in
\(\mathfrak C\).  By item~1 it is realized in \(\mathfrak C\) by some
\((b'_i)_{i\geq1}\).  Put \(b'_0=b\).  Then \((b'_i)\) has the same type
over \(A\) as \((b_i)\), so it is \(A\)-indiscernible and
\(\{\varphi(x;b'_i)\}\) is inconsistent.  If \(\psi(x;c)\) forks as
witnessed by \(b_1,\ldots,b_m\) in \(\mathcal N\) and formulas
\(\varphi_j(x;b_j)\) dividing over \(A\), realize
\(\operatorname{tp}^{\mathcal N}(b_1\cdots b_m/Ac)\) in \(\mathfrak C\) in
the same way by \(b'_1\cdots b'_m\).  The implication
\(\forall x(\psi(x;c)\to\bigvee_j\varphi_j(x;b_j))\) transfers.  Each
\(\varphi_j(x;b_j)\) divides as witnessed in some
\(\mathcal N'\succcurlyeq\mathcal N\); the type over \(Ab'_j\) obtained from
the type of that witness sequence over \(Ab_j\) by replacing \(b_j\) with
\(b'_j\) is finitely satisfiable in \(\mathcal N'\), because
\(b_j\equiv_Ab'_j\) there, and therefore in \(\mathfrak C\); realizing it
by item~1 gives a dividing witness for \(\varphi_j(x;b'_j)\) in
\(\mathfrak C\).

Now let \(\mathcal M\preccurlyeq\mathfrak C\) contain the parameters.
Witnesses in \(\mathfrak C\) are witnesses in an elementary extension of
\(\mathcal M\).  Conversely, let witnesses lie in some
\(\mathcal N\succcurlyeq\mathcal M\).  Both \(\mathcal N\) and
\(\mathfrak C\) are elementary extensions of \(\mathcal M\), so the
structures with all elements of \(\mathcal M\) named are elementarily
equivalent.  By \cref{thm:p03-w03-p03-006-elementary-amalgamation}~(1)
there are \(\mathcal D\succcurlyeq\mathfrak C\) and an elementary
embedding \(g:\mathcal N\to\mathcal D\) fixing \(\dom\mathcal M\)
pointwise.  The images under \(g\) are witnesses in \(\mathcal D\), since
the defining conditions are first-order conditions over the parameters
in \(\mathcal M\), and the first part moves them into \(\mathfrak C\).

\emph{Item 4.}  Adding unused object variables changes no consistency
question, since those variables are unconstrained.  For parameters, after
rearranging and repeating coordinates we may write \(e'=ee''\), with
\(e''\) possibly infinite and small.  Projecting an
\(A\)-indiscernible sequence of such pairs to its first coordinates shows
that dividing with the longer parameter tuple implies dividing with the
shorter one.

Conversely, let \((e_i:i<\omega)\) witness dividing, with \(e_0=e\),
and choose \(k\geq1\) so that the instances \(\varphi(x;e_i)\) are
\(k\)-inconsistent.  Item~2 gives \(e''_i\) with
\(e_ie''_i\equiv_Aee''\) for each \(i\).  Introduce variables
\((z_i,z''_i)\), \(i<\omega\), of the same coordinate lengths, and require
that the pairs be \(A\)-indiscernible, that every pair realize
\(\operatorname{tp}(ee''/A)\), and that every increasing \(k\)-tuple
of their first coordinates satisfy
\(\neg\exists x\bigwedge_{j<k}\varphi(x;z_{i_j})\).

Every finite part of these requirements is satisfiable.  It mentions only
finitely many formulas and finitely many coordinates in each block.
Colour the finite subsets of \(\omega\) by the truth values of those
formulas on the corresponding increasing tuples of \((e_ie''_i)\).
Repeated infinite Ramsey thinning makes all these finitely many
colourings constant on one infinite subset.  Assign the mentioned
indices, in order, to members of that subset.  The singleton types and
the \(k\)-inconsistency conditions already hold on the original sequence
and survive thinning.  Item~1 realizes the whole requirement, giving an
\(A\)-indiscernible sequence \((f_if''_i:i<\omega)\), with each pair of
type \(ee''\) over \(A\), and with the \(\varphi(x;f_i)\)
\(k\)-inconsistent.  Item~2 copies this sequence so that its first pair is
exactly \(ee''\).  This witnesses dividing in the longer presentation.
The argument uses only small tuples and finitely many coordinates in each
Ramsey step, so it covers infinite parameter padding as well as finite
padding.

Forking is unchanged too: the implication witnessing it and the dividing
disjuncts are the same formulas under either presentation.  The same
observation gives the assertion about unused object variables for forking.

\end{proof}
\end{annexitem}

\begin{annexitem}{annex:p03-w03-p03-031-ext}{thm:p03-w03-p03-031-extraction}{Indiscernible extraction copying finite types}
\begin{proof}
Write \(\beth_\alpha\) for \(\beth_\alpha(\lambda)\) and
\(\kappa=\beth_\mu\).  For a set \(s\) of \(n+1\) indices with increasing
enumeration \(s_0<\cdots<s_n\), let \(\chi_n(s)\) be the complete type
of \(c_{s_0}\cdots c_{s_n}\) over \(B\).  There are at most
\(2^\lambda=\beth_1\) complete types over \(B\) in \((n+1)n_0\) variables,
since there are at most \(\lambda\) formulas over \(B\).  For
\(1\leq\alpha<\mu\) we have \(2^\lambda\leq\beth_\alpha\), and for a
finite \(n\) the finite beth iteration of
\cref{thm:p06-erdos-rado} satisfies
\(\beth_n(\beth_\alpha)=\beth_{\alpha+n}\).  So
\cref{thm:p06-erdos-rado} gives
\[
 \beth_{\alpha+n}^+\longrightarrow(\beth_\alpha^+)^{n+1}_{\beth_\alpha}.
\]
A partition relation for a cardinal applies to every set of that
cardinality through a bijection, since it concerns subsets only.  The
cardinal \(\mu=(2^\lambda)^+\) is a successor, hence regular, and exceeds
\(2^\lambda\).

We build, by recursion on \(n\), complete types \(p_n\) over \(B\) in
\((n+1)n_0\) variables, sets \(S_n\subseteq[1,\mu)\) of cardinality
\(\mu\), and sets \(H_{n,\alpha}\subseteq\kappa\) for \(\alpha\in S_n\) with
\(\abs{H_{n,\alpha}}=\beth_\alpha^+\), such that for every \(m\leq n\) every
increasing \((m+1)\)-tuple of indices from \(H_{n,\alpha}\) gives a tuple
of \(c\)'s realizing \(p_m\).

For \(n=0\) and each \(\alpha\in[1,\mu)\): the \(\kappa\) indices fall into
at most \(2^\lambda\leq\beth_\alpha\) classes according to
\(\chi_0\).  If every class had at most \(\beth_\alpha\) members, then
\(\kappa\leq\beth_\alpha\cdot\beth_\alpha=\beth_\alpha\), which is false.
Let \(H_{0,\alpha}\) be a subset of size \(\beth_\alpha^+\) of a larger
class, and \(q_{0,\alpha}\) its type.  The map
\(\alpha\mapsto q_{0,\alpha}\) on \([1,\mu)\) takes at most \(2^\lambda<\mu\)
values; by regularity of \(\mu\) one value \(p_0\) is taken on a set
\(S_0\) of cardinality \(\mu\).

Given stage \(n\), let \(\alpha\in[1,\mu)\).  Since \(S_n\) has cardinality
\(\mu\) and \(\mu\) is regular, \(S_n\) is unbounded in \(\mu\); pick
\(\beta\in S_n\) with \(\beta\geq\alpha+n+1\).  Then
\(\abs{H_{n,\beta}}=\beth_\beta^+\geq\beth_{\alpha+n+1}^+\).  The relation
\(\beth_{\alpha+n+1}^+\to(\beth_\alpha^+)^{n+2}_{\beth_\alpha}\) applied to
\(\chi_{n+1}\) on a subset of \(H_{n,\beta}\) of cardinality
\(\beth_{\alpha+n+1}^+\) gives \(H_{n+1,\alpha}\subseteq H_{n,\beta}\) of
cardinality \(\beth_\alpha^+\) on whose increasing \((n+2)\)-tuples
\(\chi_{n+1}\) is constant, with value \(q_{n+1,\alpha}\).  Increasing
\((m+1)\)-tuples from \(H_{n+1,\alpha}\) with \(m\leq n\) come from
\(H_{n,\beta}\), so they realize \(p_m\).  As before some value
\(p_{n+1}\) of \(\alpha\mapsto q_{n+1,\alpha}\) is taken on a set
\(S_{n+1}\) of cardinality \(\mu\), and the sets \(H_{n+1,\alpha}\),
\(\alpha\in S_{n+1}\), satisfy the requirement at stage \(n+1\).

Each \(H_{n,\alpha}\) is infinite.  Hence for \(m\leq n\) the restriction of
\(p_n\) to any increasing choice of \(m+1\) of its \(n+1\) blocks of
variables is \(p_m\), because such subtuples of a realization taken from
\(H_{n,\alpha}\) realize \(p_m\).  Let \(\Sigma\) be the set of formulas, in
variables \((x_i)_{i<\omega}\) of length \(n_0\), saying that for every
\(n\) and every \(i_0<\cdots<i_n\) the tuple
\(x_{i_0}\cdots x_{i_n}\) realizes \(p_n\).  A finite part of
\(\Sigma\) uses indices below some \(N\geq1\).  Take \(\alpha\in S_{N-1}\)
and the first \(N\) members \(h_0<\cdots<h_{N-1}\) of \(H_{N-1,\alpha}\),
and put \(x_i=c_{h_i}\).  Every increasing \((m+1)\)-subtuple with
\(m\leq N-1\) realizes \(p_m\), so the finite part is satisfied.
By \cref{thm:p03-w03-p03-031-saturation}~(1), \(\Sigma\) is realized by
some \((b_i)_{i<\omega}\).  All increasing \((n+1)\)-tuples of it realize
the complete type \(p_n\) over \(B\), so the sequence is
\(B\)-indiscernible, and \(b_0\cdots b_n\) realizes \(p_n\), which is the
type over \(B\) of \(c_{h_0}\cdots c_{h_n}\) for any increasing
\(h_0<\cdots<h_n\) from \(H_{n,\alpha}\), \(\alpha\in S_n\).
\end{proof}
\end{annexitem}

\begin{annexitem}{annex:p03-w03-p03-031-ndi}{thm:p03-w03-p03-031-nondividing-indiscernible}{Nondividing and indiscernible sequences}
\begin{proof}
We give the argument of \cite[Lemma~3.1, p.~267]{Casanovas2007Simplicity}
with the small tuples and arbitrary infinite index orders allowed in the
statement.  Let \(I=(b_i:i\in X)\) be \(A\)-indiscernible, with a first
index \(0\) and \(b_0=b\).  A formula in \(\operatorname{tp}(a/Ab)\)
uses finitely many coordinates of \(a,b\) and finitely many parameters
\(d\) from \(A\); write it as \(\varphi(x;b,d)\).  By
\cref{thm:p03-w03-p03-031-saturation}\,(4), the unused coordinates cause
no change to dividing.  In any \(A\)-indiscernible sequence beginning
with \(bd\), the coordinates for \(d\) remain equal to \(d\).

\emph{(1) implies (2).}  Put \(p(x,b)=\operatorname{tp}(a/Ab)\) and
\(q(x)=\bigcup_{i\in X}p(x,b_i)\).  If \(q\) were inconsistent,
conjoining the finitely many formula schemes in a finite inconsistent
subset would give \(\varphi(x;b,d)\in p\) and finitely many indices
\(i_1<\cdots<i_k\) with
\(\neg\exists x\bigwedge_{j=1}^k\varphi(x;b_{i_j},d)\).
Indiscernibility makes this true on every increasing \(k\)-tuple from
\(I\).  Compactness realizes the same increasing finite-tuple types on
an \(\omega\)-indexed sequence: every finite requirement can be read
on sufficiently many increasing indices of the infinite order \(X\).
Copy its first member to \(b\) by
\cref{thm:p03-w03-p03-031-saturation}\,(2).  This is a dividing witness
for \(\varphi(x;b,d)\), contradicting (1).  Therefore \(q\) is
consistent; let \(c\) realize it by saturation.

Let \(\Sigma(x)\) express that \(I\) is indiscernible over \(Ax\):
for every formula \(\psi\) with parameters from \(A\) and every pair of
increasing \(n\)-tuples of indices, include the equality of the truth
values of \(\psi(x,b_{i_1},\ldots,b_{i_n})\) and
\(\psi(x,b_{j_1},\ldots,b_{j_n})\).  We verify finite satisfiability of
\(q\cup\Sigma\).  A finite part involves finitely many indices
\(i_1<\cdots<i_m\) and finitely many formula schemes.  On a countably
infinite subset of \(X\), colour each finite subset of the relevant
size by the value of its formula at \(c\), using its increasing order
in \(X\).  The infinite Ramsey theorem, applied finitely many times,
gives an infinite subset homogeneous for every mentioned scheme.  Choose
\(j_1<\cdots<j_m\) in it.  The corresponding instances of \(q\) hold
at \(c\), and all the required indiscernibility equalities hold there.
Since \(b_{j_1}\cdots b_{j_m}\equiv_A b_{i_1}\cdots b_{i_m}\),
\cref{thm:p03-w03-p03-031-saturation}\,(2) supplies \(c'\) copying the
whole configuration to the original indices.  It realizes the fixed
finite part of \(q\cup\Sigma\).  This argument needs neither a
subsequence of order type \(X\) nor finite coordinate blocks.

Saturation now gives \(a'\) realizing \(q\cup\Sigma\).  Since the
instance of \(p\) at index zero is included, \(a'\equiv_{Ab}a\); and
\(\Sigma\) says exactly that \(I\) is \(Aa'\)-indiscernible.

\emph{(2) implies (3).}  Copy the configuration \(a'I\) over \(Ab\)
so that \(a'\) is sent to \(a\), using
\cref{thm:p03-w03-p03-031-saturation}\,(2).  Its sequence part \(J\)
satisfies \(a'I\equiv_{Ab}aJ\).  Thus \(J\equiv_{Ab}I\), its first
member is \(b\), and it is \(Aa\)-indiscernible.

\emph{(3) implies (1).}  It is enough to test each
\(\varphi(x;b,d)\in\operatorname{tp}(a/Ab)\) on an
\(\omega\)-indexed dividing candidate \(I\) beginning with \(b\).
Take \(J=(b'_i:i<\omega)\equiv_{Ab}I\) that is
\(Aa\)-indiscernible.  Since \(b'_0=b\),
\(\varphi(a;b'_0,d)\) holds; indiscernibility over \(Aa\) gives
\(\varphi(a;b'_i,d)\) for every \(i\).  The instances along \(J\)
are therefore consistent.  Equality of the types of \(J\) and \(I\)
over \(A\) transfers every finite consistency assertion, so the
instances along \(I\) are consistent too.  No dividing candidate
witnesses dividing, proving (1).
\end{proof}
\end{annexitem}

\begin{annexitem}{annex:p03-w03-p03-031-pairs}{thm:p03-w03-p03-031-pairs}{Nondividing of pairs and enlargement of the base}
\begin{proof}
\emph{Item 1.}  This is \cite[Proposition~3.4, p.~267]{Casanovas2007Simplicity},
which Casanovas leaves as an application of his Lemma~3.1.  Let
\(\psi(x,y;c)\in\operatorname{tp}(ab/B)\) with \(c\) finite from \(B\), and
let \(I\) be \(A\)-indiscernible with first member \(c\).  The formulas of
\(\operatorname{tp}(a/Ac)\) belong to \(\operatorname{tp}(a/B)\), so
\(\operatorname{tp}(a/Ac)\) does not divide over \(A\).  By
\cref{thm:p03-w03-p03-031-nondividing-indiscernible}, (1)\(\Rightarrow\)(3),
there is \(J\equiv_{Ac}I\) that is \(Aa\)-indiscernible; it begins with
\(c\).  Likewise \(\operatorname{tp}(b/Aac)\subseteq
\operatorname{tp}(b/Ba)\) does not divide over \(Aa\).  Applying
(1)\(\Rightarrow\)(2) of the same lemma over the base \(Aa\) to the
\(Aa\)-indiscernible \(J\) gives \(b'\equiv_{Aac}b\) with \(J\)
\(Aab'\)-indiscernible.  From \(\psi(a,b;c)\) we get \(\psi(a,b';c)\) and
then \(\psi(a,b';c')\) for every member \(c'\) of \(J\).  So
\(\{\psi(x,y;c'):c'\in J\}\) is consistent, and since \(J\equiv_AI\) so is
\(\{\psi(x,y;c'):c'\in I\}\).  Hence \(\psi(x,y;c)\) does not divide over
\(A\).

\emph{Item 2.}  This is \cite[Proposition~3.5, p.~267]{Casanovas2007Simplicity}.
Let \(I\) be an \(A\)-indiscernible sequence with first member \(a\)
such that the set \(\{\varphi(x;a'):a'\in I\}\) is inconsistent.  Applying
\cref{thm:p03-w03-p03-031-nondividing-indiscernible}, (1)\(\Rightarrow\)(3),
with \(b\) as realizing tuple and \(a\) as parameter tuple, gives
\(J\equiv_{Aa}I\) that is \(Ab\)-indiscernible.  It begins with \(a\), and
\(\{\varphi(x;a'):a'\in J\}\) is inconsistent.  So \(\varphi(x;a)\) divides
over \(Ab\).
\end{proof}
\end{annexitem}

\begin{annexitem}{annex:p03-w03-p03-031-basic}{thm:p03-w03-p03-031-basic-calculus}{Invariance, monotonicity and extension in every complete theory}
\begin{proof}
\emph{Item 1.}  Let \(I=(b_i)_{i<\omega}\) witness that \(\varphi(x;b)\)
divides over \(A\), with \(b_0=b\).  By
\cref{thm:p03-w03-p03-031-saturation}~(2) there is \((b'_i)_{i\geq1}\) with
\(b(b_i)_{i\geq1}\equiv_Ab'(b'_i)_{i\geq1}\).  The sequence
\((b',b'_1,\ldots)\) has the type of \(I\) over \(A\), so it is
\(A\)-indiscernible, begins with \(b'\), and makes
\(\{\varphi(x;b'_i)\}\) inconsistent.  The converse is symmetric.  If
\(\psi(x;b)\) implies \(\bigvee_j\varphi_j(x;e_j)\) with each
\(\varphi_j(x;e_j)\) dividing over \(A\), take \(e'\) with
\(be\equiv_Ab'e'\); the implication is a formula over \(A\) in \(be\), so
it holds for \(b'e'\), and each \(\varphi_j(x;e'_j)\) divides by the
dividing case.  For the statement about types, an elementary map \(f\)
fixing \(A\) satisfies \(f(e)\equiv_Ae\) for every tuple \(e\) from its
domain, so \(\varphi(x;e)\in p\) forks over \(A\) exactly when
\(\varphi(x;f(e))\in f(p)\) does.

\emph{Item 2.}  A dividing formula implies itself, a one-member
disjunction.  If \(\psi(x;c)\) implies \(\bigvee_{j\leq n}\theta_j\) and
each \(\theta_j\) implies a finite disjunction of formulas dividing over
\(A\), then \(\psi(x;c)\) implies the finite disjunction of all of these.

\emph{Item 3.}  The formulas of \(p\mathbin{\upharpoonright}B\) are formulas
of \(p\), so a forking formula in the former would lie in the latter.  If
\(\varphi(x;e)\) divides over \(B\), its witness sequence is
\(B\)-indiscernible, hence \(A\)-indiscernible, so \(\varphi(x;e)\) divides
over \(A\); applying this to the disjuncts shows that a formula forking over
\(B\) forks over \(A\).  Hence if \(p\) forked over \(B\) it would contain a
formula forking over \(A\).  For the tuple \(ab\), a formula
\(\varphi(x;e)\in\operatorname{tp}(a/C)\) is, with the variables \(y\) of
\(b\) added unused, a formula of \(\operatorname{tp}(ab/C)\), and by
\cref{thm:p03-w03-p03-031-saturation}~(4) it forks, or divides, in one
reading if and only if in the other.  The dividing versions are proved by
the same three arguments with ``divides'' in place of ``forks''.

\emph{Item 4.}  A forking formula of a complete type over \(C\) has finitely
many parameters \(c\) from \(C\) and lies in
\(\operatorname{tp}(a/Ac)\) for any realization \(a\); conversely a forking
formula of \(\operatorname{tp}(a/Ac)\), with \(c\) from \(C\), lies in the
type over \(C\) that \(a\) realizes.

\emph{Item 5.}  Let \(\Gamma=\pi\cup\{\neg\varphi(x;e):e\text{ from }C,\
\varphi(x;e)\text{ forks over }A\}\).  If \(\Gamma\) were not finitely
satisfiable, some finite conjunction \(\psi\) of members of \(\pi\) would
imply a finite disjunction of formulas forking over \(A\), and by item~2
\(\psi\) would fork over \(A\), contrary to the hypothesis on \(\pi\).  By
\cref{thm:p03-w03-p03-031-saturation}~(1) some \(a\) realizes \(\Gamma\).
Then \(\operatorname{tp}(a/C)\) contains \(\pi\), and it contains no
formula \(\varphi(x;e)\) forking over \(A\), since it contains
\(\neg\varphi(x;e)\).  For a complete type over \(B\) take \(\pi\) to be
that type.
\end{proof}
\end{annexitem}

\begin{annexitem}{annex:p03-w03-p03-031-simp}{thm:p03-w03-p03-031-simple-extension}{Existence and extension in simple theories}
\begin{proof}
By \cref{thm:p03-w03-p03-031-saturation}~(3), forking in the sense of
\cref{def:p03-dividing-forking} for the model \(\mathfrak C\) is forking
witnessed inside \(\mathfrak C\), so \cref{def:p03-simple-theory}, applied
to the model \(\mathfrak C\), gives local character for the forking
relation used here.

\emph{Item 3.}  For a complete type \(p\) over \(C\) in a finite tuple
\(x\) of variables, \cref{def:p03-simple-theory} gives \(A_0\subseteq C\)
with \(\abs{A_0}\leq\abs T+\abs x=\abs T\), since \(\abs T\) is infinite,
such that \(p\) does not fork over \(A_0\).  If
\(A_0\subseteq A'\subseteq C\), monotonicity in the base
(\cref{thm:p03-w03-p03-031-basic-calculus}~(3)) shows that \(p\) does not
fork over \(A'\).

\emph{Item 1.}  Apply item~3 with \(C=A\) and \(A'=A\).  For a type in an
infinite tuple of variables, a forking formula would involve finitely many
of them; its restriction to those variables is a type in finitely many
variables that does not fork over \(A\), and
\cref{thm:p03-w03-p03-031-saturation}~(4) shows that the forking formula
would fork as a formula in those variables.

\emph{Item 2.}  By item~1 a complete type \(p\) over \(A\) does not fork
over \(A\); \cref{thm:p03-w03-p03-031-basic-calculus}~(5), with base \(A\),
\(B=A\) and \(C\), gives the required extension.
\end{proof}
\end{annexitem}

\begin{annexitem}{annex:p03-w03-p03-031-morley}{thm:p03-w03-p03-031-morley-existence}{Morley sequences}
\begin{proof}
\emph{Item 1.}  We follow \cite[Lemma~5.1, p.~272]{Casanovas2007Simplicity}.
Let \(n_0\) be the length of the variable tuple of \(p\),
\(\lambda=\abs T+\abs B\) and \(\kappa=\beth_{(2^\lambda)^+}(\lambda)\),
which is small by the convention.  Choose \(c_i\) for \(i<\kappa\) by
recursion: given \((c_j)_{j<i}\), the type \(p\) does not fork over \(A\),
so by \cref{thm:p03-w03-p03-031-basic-calculus}~(5) it has an extension to
\(B\cup\{c_j:j<i\}\) not forking over \(A\); let \(c_i\) realize it
(\cref{thm:p03-w03-p03-031-saturation}~(1)).  By
\cref{thm:p03-w03-p03-031-extraction}, applied over \(B\), there is a
\(B\)-indiscernible \((a_i)_{i<\omega}\) such that for every \(n\) some
\(j_0<\cdots<j_n\) satisfy \(a_0\cdots a_n\equiv_Bc_{j_0}\cdots c_{j_n}\).
Then \(a_n\equiv_Bc_{j_n}\), so \(a_n\) realizes \(p\).  The sequence is
\(A\)-indiscernible because \(A\subseteq B\).  The type
\(\operatorname{tp}(c_{j_n}/B\cup\{c_{j_0},\ldots,c_{j_{n-1}}\})\) is a
restriction of \(\operatorname{tp}(c_{j_n}/B\cup\{c_j:j<j_n\})\), so it does
not fork over \(A\) (\cref{thm:p03-w03-p03-031-basic-calculus}~(3)).  The
map fixing \(B\) and sending \(c_{j_0}\cdots c_{j_n}\) to \(a_0\cdots a_n\)
is elementary and carries this type to
\(\operatorname{tp}(a_n/B\cup\{a_0,\ldots,a_{n-1}\})\); by invariance
(\cref{thm:p03-w03-p03-031-basic-calculus}~(1)) the latter does not fork
over \(A\), and neither does its restriction
\(\operatorname{tp}(a_n/A\cup\{a_0,\ldots,a_{n-1}\})\).  So
\((a_i)_{i<\omega}\) is a Morley sequence over \(A\) in \(p\) and is
\(B\)-indiscernible.

\emph{Item 2.}  We follow \cite[Remark~5.2, p.~272]{Casanovas2007Simplicity}.
A dividing formula of
\(\operatorname{tp}((a_i)_{i\in Z}/A\cup\{a_i:i\in Y\})\) uses variables
for finitely many \(i\in Z\), and by
\cref{thm:p03-w03-p03-031-saturation}~(4) it divides as a formula in those
variables; so it suffices to treat finite \(Z\), by induction on
\(\abs Z\).  Write \(a_Y\) for \((a_i)_{i\in Y}\).  For \(Z=\varnothing\)
the type consists of true sentences over \(Aa_Y\), and a consistent formula
without free variables does not divide.  Let \(Z=Z'\cup\{z\}\) with \(z\)
the largest element of \(Z\), and assume
\(\operatorname{tp}(a_{Z'}/Aa_Y)\) does not divide over \(A\).  All indices
in \(Y\cup Z'\) are below \(z\), so
\(\operatorname{tp}(a_z/Aa_Ya_{Z'})\) is a restriction of
\(\operatorname{tp}(a_z/A\cup\{a_j:j<z\})\), which does not fork over
\(A\).  Hence it does not fork, and so does not divide, over \(A\)
(\cref{thm:p03-w03-p03-031-basic-calculus}~(2),~(3)), and by monotonicity
in the base it does not divide over \(Aa_{Z'}\).
\Cref{thm:p03-w03-p03-031-pairs}~(1), with \(Aa_Y\) as the larger set,
\(a_{Z'}\) as the first tuple and \(a_z\) as the second, shows that
\(\operatorname{tp}(a_{Z'}a_z/Aa_Y)\) does not divide over \(A\).
\end{proof}
\end{annexitem}

\begin{annexitem}{annex:p03-w03-p03-031-kim}{thm:p03-w03-p03-031-kim-lemma}{Kim's lemma}
\begin{proof}
We follow \cite[Proposition~5.3, p.~272]{Casanovas2007Simplicity}.

\emph{From inconsistency to dividing.}  By
\cref{thm:p03-w03-p03-031-simple-extension}~(1),
\(\operatorname{tp}(a/A)\) does not fork over \(A\), and
\cref{thm:p03-w03-p03-031-morley-existence}~(1), with \(B=A\), gives an
infinite Morley sequence \((a_i)\) over \(A\) in it.  By hypothesis
\(\{\varphi(x;a_i)\}\) is inconsistent, so \(\varphi(x;a_0)\) divides over
\(A\), and so does \(\varphi(x;a)\) by invariance, since \(a_0\equiv_Aa\).

\emph{From dividing to inconsistency.}  Suppose \(\varphi(x;a)\) divides over
\(A\) but \((a_i)_{i<\omega}\) is a Morley sequence over \(A\) in
\(\operatorname{tp}(a/A)\) with \(\{\varphi(x;a_i)\}\) consistent.  Let
\(X\) be the set \(\abs T^+\) ordered by the reverse of the ordinal order;
write \(<_X\).  The set of formulas in variables \((y_i)_{i\in X}\) saying
that every \(<_X\)-increasing \((n+1)\)-tuple realizes
\(\operatorname{tp}(a_0\cdots a_n/A)\), for every \(n\), is finitely
satisfied by initial pieces of \((a_i)_{i<\omega}\), so by
\cref{thm:p03-w03-p03-031-saturation}~(1) it is realized by some
\((a'_i)_{i\in X}\).  This sequence is \(A\)-indiscernible, each
\(a'_i\equiv_Aa\), and every finite subset of
\(\{\varphi(x;a'_i):i\in X\}\) is satisfiable because the corresponding
existential formula holds of \(a_0\cdots a_n\).  It is \(A\)-independent:
if \(\operatorname{tp}(a'_i/A\cup\{a'_j:j<_Xi\})\) forked over \(A\), then
by finite character some \(j_1<_X\cdots<_Xj_m<_Xi\) would make
\(\operatorname{tp}(a'_i/Aa'_{j_1}\cdots a'_{j_m})\) fork over \(A\); but
\(a'_{j_1}\cdots a'_{j_m}a'_i\equiv_Aa_0\cdots a_m\), and
\(\operatorname{tp}(a_m/Aa_0\cdots a_{m-1})\) does not fork over \(A\), so
invariance gives a contradiction.

Let \(c\) realize \(\{\varphi(x;a'_i):i\in X\}\)
(\cref{thm:p03-w03-p03-031-saturation}~(1)) and put
\(D=A\cup\{a'_i:i\in X\}\).  By
\cref{thm:p03-w03-p03-031-simple-extension}~(3) there is \(A_0\subseteq D\)
with \(\abs{A_0}\leq\abs T\) such that \(\operatorname{tp}(c/D)\) does not
fork over \(A_0\).  Choose for each element of \(A_0\setminus A\) one index
\(i\) such that it is an entry of \(a'_i\), and let \(Y\) be the set of the
chosen indices; \(\abs Y\leq\abs T\).  Then
\(A_0\subseteq A\cup a'_Y\subseteq D\), so
\(\operatorname{tp}(c/D)\) does not fork over \(Aa'_Y\).  As \(\abs T^+\) is
regular and \(\abs Y\leq\abs T\), the ordinal
\(\iota=\sup\{y+1:y\in Y\}\) is below \(\abs T^+\), and \(\iota<_Xy\) for
every \(y\in Y\).  By \cref{thm:p03-w03-p03-031-morley-existence}~(2)
\(\operatorname{tp}(a'_Y/Aa'_\iota)\) does not divide over \(A\).  The
formula \(\varphi(x;a'_\iota)\) divides over \(A\), by invariance, so
\cref{thm:p03-w03-p03-031-pairs}~(2) shows that it divides over
\(Aa'_Y\).  It belongs to \(\operatorname{tp}(c/D)\) and forks over
\(Aa'_Y\) (\cref{thm:p03-w03-p03-031-basic-calculus}~(2)), which
contradicts the choice of \(Y\).
\end{proof}
\end{annexitem}

\begin{annexitem}{annex:p03-w03-p03-031-fd}{thm:p03-w03-p03-031-forking-dividing}{Forking equals dividing in simple theories}
\begin{proof}
We follow \cite[Proposition~5.4, p.~272]{Casanovas2007Simplicity}.  Both
notions are defined for partial types through finite conjunctions, so it
suffices to treat a formula \(\psi(x;a)\) with \(a\) finite.  Dividing
implies forking by \cref{thm:p03-w03-p03-031-basic-calculus}~(2).
Suppose \(\psi(x;a)\) does not divide over \(A\) but implies
\(\bigvee_{k=1}^n\varphi_k(x;a_k)\) with each \(\varphi_k(x;a_k)\)
dividing over \(A\).  By \cref{thm:p03-w03-p03-031-simple-extension}~(1)
and \cref{thm:p03-w03-p03-031-morley-existence}~(1) there is a Morley
sequence \((a^ja^j_1\cdots a^j_n)_{j<\omega}\) over \(A\) in
\(\operatorname{tp}(aa_1\cdots a_n/A)\).  The sequence \((a^j)_j\) is
\(A\)-indiscernible with \(a^0\equiv_Aa\); by invariance
\(\psi(x;a^0)\) does not divide over \(A\), so
\(\{\psi(x;a^j):j<\omega\}\) is realized by some \(c\).  The implication
is a formula over \(A\) true of \(aa_1\cdots a_n\), so it holds for each
\(a^ja^j_1\cdots a^j_n\), and \(c\) satisfies some
\(\varphi_{k_j}(x;a^j_{k_j})\) for each \(j\).  Some \(k\) equals
\(k_j\) for infinitely many \(j\); call that infinite set \(X\).  For any
finite \(F\subseteq\omega\), choose \(F'\subseteq X\) of the same size;
indiscernibility transfers the satisfiability of
\(\{\varphi_k(x;a^j_k):j\in F'\}\) to
\(\{\varphi_k(x;a^j_k):j\in F\}\).  So
\(\{\varphi_k(x;a^j_k):j<\omega\}\) is consistent.  The sequence
\((a^j_k)_{j<\omega}\) is \(A\)-indiscernible, each term realizes
\(\operatorname{tp}(a_k/A)\), and
\(\operatorname{tp}(a^j_k/A\cup\{a^l_k:l<j\})\) is obtained from the
nonforking type
\(\operatorname{tp}(a^ja^j_1\cdots a^j_n/A\cup\{a^la^l_1\cdots a^l_n:l<j\})\)
by restricting the realizing tuple and the parameters, so it does not fork
over \(A\) (\cref{thm:p03-w03-p03-031-basic-calculus}~(3)).  Thus it is a
Morley sequence over \(A\) in \(\operatorname{tp}(a_k/A)\), and
\cref{thm:p03-w03-p03-031-kim-lemma} shows that \(\varphi_k(x;a_k)\) does
not divide over \(A\), a contradiction.  For a complete type the last
sentence of the theorem follows, since it forks exactly when it contains a
forking formula.
\end{proof}
\end{annexitem}

\begin{annexitem}{annex:p03-w03-p03-031-sym}{thm:p03-w03-p03-031-symmetry}{Symmetry}
\begin{proof}
We follow \cite[Proposition~5.5, p.~273]{Casanovas2007Simplicity}.

\emph{Reduction to tuples.}  Assume
\(A\mathrel{\raisebox{-0.3ex}{\(\smile\)}\mkern-11.5mu\mid\mkern6.5mu}_{C}B\)
and let \(b\) be finite from \(B\).  If \(\operatorname{tp}(b/A\cup C)\)
forked over \(C\), it would contain a forking formula with finitely many
parameters \(a\) from \(A\), and \(\operatorname{tp}(b/Ca)\) would fork over
\(C\).  But \(\operatorname{tp}(a/Cb)\) is a restriction of
\(\operatorname{tp}(a/B\cup C)\), which does not fork over \(C\).  So it is
enough to show for finite \(a,b\): if \(\operatorname{tp}(a/Cb)\) does not
fork over \(C\), then neither does \(\operatorname{tp}(b/Ca)\); the reverse
implication of the tuple statement is the same claim with \(a\) and \(b\)
exchanged.  By \cref{thm:p03-w03-p03-031-forking-dividing} it suffices to
show that \(\operatorname{tp}(b/Ca)\) does not divide over \(C\).

\emph{A Morley sequence through \(a\).}  By
\cref{thm:p03-w03-p03-031-morley-existence}~(1), with base \(C\) and
domain \(Cb\), there is a Morley sequence \((a_i)_{i<\omega}\) over \(C\)
in \(\operatorname{tp}(a/Cb)\) that is \(Cb\)-indiscernible.  Since
\(a_0\equiv_{Cb}a\), \cref{thm:p03-w03-p03-031-saturation}~(2) gives a
sequence with the same type over \(Cb\) whose first member is \(a\); it
has the same three properties, being independent by invariance.  So we may
assume \(a_0=a\).

\emph{Adding base parameters.}  If \(d\) is from \(C\) and
\(\operatorname{tp}(a_i/C\cup\{a_j:j<i\})\) does not fork over \(C\), then
neither does \(\operatorname{tp}(a_id/C\cup\{a_jd:j<i\})\).  Indeed,
suppose a formula \(\chi(x,z;e)\) of the latter forks over \(C\) as a
formula in \(x,z\), say \(\chi(x,z;e)\to\bigvee_k\theta_k(x,z;e_k)\) with
each \(\theta_k\) dividing over \(C\) by a sequence \((e_k^l)_l\).  Then
\(\chi(x,d;e)\) belongs to \(\operatorname{tp}(a_i/C\cup\{a_j:j<i\})\) and
implies \(\bigvee_k\theta_k(x,d;e_k)\).  Each \(\theta_k(x,d;e_k)\), with
parameter tuple \(de_k\), divides over \(C\): the sequence \((de_k^l)_l\)
is \(C\)-indiscernible because \(d\) is from \(C\), and
\(\{\theta_k(x,d;e_k^l)\}\) is inconsistent because
\(\{\theta_k(x,z;e_k^l)\}\) is.  So \(\chi(x,d;e)\) forks over \(C\), a
contradiction.  Hence for every finite \(d\) from \(C\) the sequence
\((a_id)_{i<\omega}\) is a Morley sequence over \(C\) in
\(\operatorname{tp}(ad/C)\).

\emph{Conclusion.}  Let \(\theta(y;a,d)\) be a formula of
\(\operatorname{tp}(b/Ca)\) with \(d\) finite from \(C\); by \cref{thm:p03-w03-p03-031-saturation}~(4) every
formula of \(\operatorname{tp}(b/Ca)\) can be presented this way.  Since
\(\theta(b;a,d)\) holds and \((a_i)\) is \(Cb\)-indiscernible,
\(\theta(b;a_i,d)\) holds for all \(i\), so
\(\{\theta(y;a_i,d):i<\omega\}\) is consistent.  As \((a_id)\) is a Morley
sequence over \(C\) in \(\operatorname{tp}(ad/C)\), Kim's lemma
(\cref{thm:p03-w03-p03-031-kim-lemma}) shows that \(\theta(y;a,d)\) does
not divide over \(C\).
\end{proof}
\end{annexitem}

\begin{annexitem}{annex:p03-w03-p03-031-trans}{thm:p03-w03-p03-031-transitivity}{Transitivity}
\begin{proof}
\emph{Item 1.}  Casanovas derives this from symmetry, his
Proposition~3.4 and the equality of forking and dividing
\cite[Proposition~5.6, p.~273]{Casanovas2007Simplicity}; we give the
derivation.  Let \(a\) be finite from \(A\); we must show that
\(\operatorname{tp}(a/D)\) does not fork over \(B\).  By symmetry
(\cref{thm:p03-w03-p03-031-symmetry}) it suffices to show that
\(\operatorname{tp}(d/Ba)\) does not fork over \(B\) for every finite \(d\)
from \(D\), and by \cref{thm:p03-w03-p03-031-forking-dividing} that it does
not divide over \(B\).  Let \(c\) enumerate \(C\).

From
\(A\mathrel{\raisebox{-0.3ex}{\(\smile\)}\mkern-11.5mu\mid\mkern6.5mu}_{B}C\)
and symmetry, every finite subtuple of \(c\) has a type over \(Ba\) not
forking over \(B\).  A dividing formula of \(\operatorname{tp}(c/Ba)\)
would use finitely many variables of \(c\) and, by
\cref{thm:p03-w03-p03-031-saturation}~(4), divide as a formula in those;
so \(\operatorname{tp}(c/Ba)\) does not divide over \(B\).
From
\(A\mathrel{\raisebox{-0.3ex}{\(\smile\)}\mkern-11.5mu\mid\mkern6.5mu}_{C}D\),
\(\operatorname{tp}(a/Cd)\) does not fork over \(C\), so by symmetry
\(\operatorname{tp}(d/Ca)\) does not fork, hence does not divide, over
\(C\).  Since \(B\subseteq C\), the set \(Ca\) is \(Ba\cup c\) and \(C\) is
\(B\cup c\).  \Cref{thm:p03-w03-p03-031-pairs}~(1), with \(B\) as the base,
\(Ba\) as the larger set, \(c\) as the first tuple and \(d\) as the second,
shows that \(\operatorname{tp}(cd/Ba)\) does not divide over \(B\).  By
monotonicity (\cref{thm:p03-w03-p03-031-basic-calculus}~(3), dividing
version), \(\operatorname{tp}(d/Ba)\) does not divide over \(B\).

\emph{Item 2.}  If \(p\) does not fork over \(A\), then it does not fork
over \(B\) and \(p\mathbin{\upharpoonright}B\) does not fork over \(A\), by
\cref{thm:p03-w03-p03-031-basic-calculus}~(3); this uses no hypothesis on
\(T\).  Conversely, let \(a\) realize \(p\).  Nonforking of
\(p\mathbin{\upharpoonright}B=\operatorname{tp}(a/B)\) over \(A\) says
\(a\mathrel{\raisebox{-0.3ex}{\(\smile\)}\mkern-11.5mu\mid\mkern6.5mu}_{A}B\),
and nonforking of \(p\) over \(B\) says
\(a\mathrel{\raisebox{-0.3ex}{\(\smile\)}\mkern-11.5mu\mid\mkern6.5mu}_{B}C\).
Item~1 with \(A\subseteq B\subseteq C\) in place of \(B\subseteq C\subseteq
D\) gives
\(a\mathrel{\raisebox{-0.3ex}{\(\smile\)}\mkern-11.5mu\mid\mkern6.5mu}_{A}C\),
that is, \(p\) does not fork over \(A\).  If \(p\) is in infinitely many
variables, apply this to the restriction of \(p\) to each finite set of
variables; a forking formula involves only finitely many.
\end{proof}
\end{annexitem}

\begin{annexitem}{annex:p03-w03-p03-031-stable}{thm:p03-w03-p03-031-stable-calculus}{The forking calculus in stable theories}
\begin{proof}
\emph{Stability gives simplicity.}  Let \(T\) be stable.  By
\cref{thm:p03-w03-p03-035-stable-simple}, \(T\) is simple in the sense of
\cref{def:p03-simple-theory}: for every model of \(T\), every parameter set
of it and every complete type over that set there is a base of size at
most \(\abs T+\abs x\) over which the type does not fork.  That proof uses,
from this chapter, only \cref{def:p03-w03-p03-031-monster} and
\cref{thm:p03-w03-p03-031-saturation}~(3) and~(4), which hold in every
complete theory; it uses none of the simple-theory results of
\cref{sec:p03-w03-p03-031-simple}, so there is no circularity in applying
those results now.  From here on \(T\) is simple, and every theorem of this
chapter whose hypothesis is ``\(T\) simple'' applies to \(T\).

\emph{Item 1.}  The first assertion and the passage to every \(A'\) with
\(A_0\subseteq A'\subseteq C\) are
\cref{thm:p03-w03-p03-031-simple-extension}~(3).  If \(T\) is superstable,
\cref{def:p03-stability-spectrum} gives a finite \(A_0\subseteq C\) over
which \(p\) does not fork in the sense of \cref{def:p03-dividing-forking}
for the model \(\mathfrak C\); by \cref{thm:p03-w03-p03-031-saturation}~(3)
this is nonforking witnessed in \(\mathfrak C\), the relation used here, and
the passage to larger \(A'\) is again monotonicity in the base
(\cref{thm:p03-w03-p03-031-basic-calculus}~(3)).

\emph{Items 2 and 3} are \cref{thm:p03-w03-p03-031-basic-calculus}~(1)
and~(3), which assume nothing about \(T\).

\emph{Item 4.}  The first two assertions are
\cref{thm:p03-w03-p03-031-simple-extension}~(1) and~(2).  The third is
\cref{thm:p03-w03-p03-031-basic-calculus}~(5) applied to the complete type
over \(B\) as the partial type \(\pi\).

\emph{Item 5} is \cref{thm:p03-w03-p03-031-symmetry}, whose statement
covers both the tuple form and the set form with the same base \(C\).

\emph{Item 6} is \cref{thm:p03-w03-p03-031-transitivity}~(2).

\emph{Item 7} is \cref{thm:p03-w03-p03-031-forking-dividing}.

\emph{The laws of \cref{thm:p03-nonforking-calculus}.}  Its item~1 is the
second assertion of item~4 above, with \(B\) in place of \(C\).  Its
item~2 is the tuple form of item~5 with base \(A\) in place of \(C\).  Its
item~3 is item~6.
\end{proof}
\end{annexitem}

\begin{annexitem}{annex:p03-w03-p03-031-rg}{thm:p03-w03-p03-031-random-graph}{Forking in the random graph}
\begin{proof}
We use \cref{thm:p03-w03-p03-015-graph-types}~(1): tuples with the same
graph pattern satisfy the same formulas.  In particular, for a finite set
\(C_0\), a tuple realizes a given complete type over \(C_0\) as soon as its
concatenation with an enumeration of \(C_0\) has the same pattern as the
corresponding concatenation for a realization.  We also use the extension
axioms of \(T_{\rm RG}\) in the following form: given finitely many
distinct vertices \(W\) and a subset \(W^+\subseteq W\), there is a vertex
outside \(W\) adjacent to every member of \(W^+\) and to no member of
\(W\setminus W^+\).  Applied repeatedly, this gives, for finitely many
pairwise distinct new vertices with any prescribed adjacency among
themselves and to \(W\), a realization outside \(W\).  Write
\(a=(a^1,\ldots,a^m)\).

\emph{Claim A: if no entry of \(a\) lies in \(B\setminus C\), then
\(\operatorname{tp}(a/B)\) does not divide over \(C\).}  Let
\(\theta(x;e)\in\operatorname{tp}(a/B)\) with
\(e=(e^1,\ldots,e^r)\) from \(B\), and let \((e_i)_{i<\omega}\) be
\(C\)-indiscernible with \(e_0=e\).  Let \(\gamma(x,y)\) be the pattern
formula of the tuple \((a,e)\); every realization of \(\gamma\) satisfies
\(\theta\), so it suffices that \(\bigcup_{i<N}\{\gamma(x,e_i)\}\) is
consistent for each \(N\).  Two facts about the sequence:
\begin{enumerate}
  \item If \(e_i^l=e_j^{l'}\) for some \(i,j\), then \(e^l=e^{l'}\).  For
        \(i=j\) this is indiscernibility.  For \(i<j\), indiscernibility
        gives \(e_0^l=e_1^{l'}\) and \(e_0^l=e_2^{l'}\), so
        \(e_1^{l'}=e_2^{l'}\), hence \(e_0^{l'}=e_1^{l'}=e_0^l\); the case
        \(i>j\) is symmetric.
  \item If \(e^l\in C\), then \(e_i^l=e^l\) for all \(i\), since
        \(y^l=e^l\) is a formula over \(C\).
\end{enumerate}
If \(a^k=e^l\), then \(a^k\in B\), so \(a^k\in C\) by hypothesis.  Put
\(x^k=a^k\) for every \(k\) with \(a^k\in C\).  For such \(k\) the
requirements of \(\gamma(x,e_i)\) relating \(x^k\) to \(e_i\) are the
atomic relations between \(a^k\) and \(e\) with \(e_i\) in place of \(e\),
and they hold because \(e_i\equiv_Ce\).  The requirements of
\(\gamma(x,e_i)\) among the entries of \(e_i\) hold because \(e_i\) has the
pattern of \(e\).  The remaining entries \(a^k\notin C\) lie outside
\(B\), so they equal no entry of \(e\) and no \(a^{k'}\in C\).  For them
we need new vertices, one for each class of equal entries, distinct from
the finite set \(W\) of the fixed \(a^{k'}\in C\) and all \(e_i^l\),
\(i<N\), with the adjacency among themselves of the corresponding entries
of \(a\), adjacent to \(a^{k'}\in C\) as \(a^k\) is, and adjacent to
\(e_i^l\) exactly when \(a^k\) is adjacent to \(e^l\).  These requirements
do not conflict: if \(e_i^l=e_j^{l'}\) then \(e^l=e^{l'}\) by fact~1, and
if \(e_i^l=a^{k'}\in C\) then \(e^l=a^{k'}\) by indiscernibility over
\(C\).  The extension axioms provide the new vertices, and together with
the fixed ones they satisfy \(\bigcup_{i<N}\gamma(x,e_i)\).

\emph{Claim B: if \(\varphi(x;e)\) divides over \(C\), every realization
has an entry among the entries of \(e\) not in \(C\).}  Otherwise a
realization \(a^*\) has no entry in \((C\cup e)\setminus C\), and Claim~A
with \(B=C\cup e\) says that \(\operatorname{tp}(a^*/Ce)\), which contains
\(\varphi(x;e)\), does not divide over \(C\).

\emph{Claim C: if no entry of \(a\) lies in \(B\setminus C\), then
\(\operatorname{tp}(a/B)\) does not fork over \(C\).}  Suppose
\(\theta(x;e)\in\operatorname{tp}(a/B)\) implies
\(\bigvee_{j\leq n}\varphi_j(x;d_j)\) with each \(\varphi_j(x;d_j)\) dividing
over \(C\).  Let \(F\) be the finite set of entries of the \(d_j\) not in
\(C\).  The set of formulas \(\operatorname{tp}(a/Ce)\cup\{x^k\neq f:f\in
F,\ a^k\notin C\}\) is finitely satisfiable: a finite part mentions a finite
\(C_0\subseteq C\); keep the entries \(a^k\in C\), and choose new vertices
for the entries \(a^k\notin C\), which lie outside \(B\), distinct from
\(C_0\), \(e\) and \(F\), with the adjacency of \(a\) to \(C_0\), to \(e\)
and among its entries.  The result has the pattern of \(a\) over
\(C_0\cup e\), so it satisfies the finite part.  By
\cref{thm:p03-w03-p03-031-saturation}~(1) some \(a^*\) realizes the whole
set.  Then \(a^*\) satisfies \(\theta(x;e)\), hence some
\(\varphi_j(x;d_j)\); but each entry of \(a^*\) is in \(C\) or outside
\(F\), so it is not an entry of \(d_j\) outside \(C\), contrary to Claim~B.

\emph{Claim D: if some entry \(a^k\) equals \(b\in B\setminus C\), then
\(\operatorname{tp}(a/B)\) forks over \(C\).}  Let \(\Gamma\) be the set of
formulas in variables \((y_i)_{i\geq1}\) consisting of
\(\operatorname{tp}(b/C)\) for each \(y_i\), together with \(y_i\neq y_j\),
\(\neg E(y_i,y_j)\), \(y_i\neq b\) and \(\neg E(y_i,b)\) for \(i\neq j\).  A
finite part mentions a finite \(C_0\subseteq C\) and finitely many
\(y_i\); new vertices outside \(C_0\cup\{b\}\), pairwise distinct and
nonadjacent, nonadjacent to \(b\), and adjacent to exactly the members of
\(C_0\) adjacent to \(b\), have the pattern of \(b\) over \(C_0\) and
satisfy it.  By \cref{thm:p03-w03-p03-031-saturation}~(1) some
\((b_i)_{i\geq1}\) realizes \(\Gamma\); put \(b_0=b\).  For every finite
\(C_0\subseteq C\), every increasing tuple of the \(b_i\) together with
\(C_0\) has the same pattern: distinct, pairwise nonadjacent entries,
outside \(C_0\) since \(b\notin C\), each with the adjacency of \(b\) to
\(C_0\).  So the sequence is \(C\)-indiscernible.  The formulas
\(x^k=b_i\) are pairwise inconsistent, so \(x^k=b\), a member of
\(\operatorname{tp}(a/B)\) by \cref{thm:p03-w03-p03-031-saturation}~(4),
divides and therefore forks over \(C\).

Claims C and D give the stated equivalence.  For local character, let
\(\mathcal M\entails T_{\rm RG}\), \(B\subseteq\dom\mathcal M\) and \(p\) a
complete type over \(B\) in finitely many variables consistent with
\(\mathcal M\).  Take the ambient model of
\cref{def:p03-w03-p03-031-monster} to be a \(\bar\kappa\)-saturated
elementary extension of \(\mathcal M\) with \(\bar\kappa>\abs{\dom\mathcal
M}\) (\cref{thm:saturated-extension}); by
\cref{thm:p03-w03-p03-031-saturation}~(3) forking in the sense of
\(\mathcal M\) is forking there.  Let \(a\) realize \(p\) and let \(A_0\) be
the finite set of entries of \(a\) lying in \(B\).  No entry of \(a\) lies
in \(B\setminus A_0\), so by Claim~C \(p\) does not fork over \(A_0\), and
\(\abs{A_0}\) is finite.  For a type in infinitely many variables, each
forking formula involves finitely many of them, and the same choice of
\(A_0\) from the entries in \(B\) gives a set of size at most
\(\abs x\leq\abs T+\abs x\).  So \(T_{\rm RG}\) satisfies
\cref{def:p03-simple-theory}.
\end{proof}
\end{annexitem}
